\section{Material}

Für die Umsetzung des Projekts wurden verschiedene
Hardware- und Softwarekomponenten verwendet.

\begin{itemize}
    \item Raspberry Pi 4 Model B
    \item JOY-PI-Erweiterungsboard
    \item Zwei DHT11-Temperatur- und Luftfeuchtigkeitssensoren
    \item I\textsuperscript{2}C-LCD-Display
    \item Relaismodul zur Lüftersteuerung
    \item Lüfter
    \item Python 3 als Programmiersprache
    \item Flask zur Erstellung der Weboberfläche
    \item SQLite zur Speicherung der Messdaten
    \item Adafruit Blinka zur Hardwareansteuerung
    \item RPi.GPIO zur Steuerung der GPIO-Pins
    \item Git zur Versionsverwaltung
    \item GitHub zur Verwaltung des Quellcodes
    \item RealVNC Viewer für den Remote-Zugriff auf den Raspberry Pi
    \item Visual Studio Code als Entwicklungsumgebung
\end{itemize}

Das GitHub-Repository des Projekts befindet sich unter:

\begin{center}
\url{https://github.com/judithhoh/LF07-Taupunkt}
\end{center}